\hspace{3mm}

\centerline{\bf \Huge Домашнее задание}
\centerline{\bf \Huge Признаки Делимости}

\begin{flushright}
    \scriptsize{РВ передаёт салам четвёртой задачей}
\end{flushright}

\problem{1!} 
Делится ли $7^{2024} + 7^{2025} + 7^{2026} + 7^{2027} \ \text{на} \ 2744?$

\problem{2!} 
Может ли сумма $2024$ последовательных натуральных чисел оканчиваться той же цифрой, что и сумма следующих $2028$ чисел?

\problem{3} 
Пятизначное число называется \textit{восхитительным}, если оно само делится на $36$, его сумма цифр делится на $36$ и его произведение цифр делится на $36$. Найдите все восхитительные числа.

\problem{4} 
Федор записал на одной доске 10000000 чисел. На второй доске Дамир записал разность каждого числа Фёдора и произведения оставшихся его чисел. Оказалось, что каждое число Дамира делится на 10000000. Докажите, что сумма квадратов чисел Фёдора делится на 10000000.